\documentclass[11pt]{article}
\usepackage[T1]{fontenc}
\usepackage{textcomp}
\usepackage[margin=0.75in]{geometry}
\usepackage{amsmath,amssymb,amsthm}
\usepackage{array,booktabs}
\usepackage{hyperref}
\setlength{\parindent}{0pt}
\newtheorem{theorem}{Theorem}
\theoremstyle{definition}
\newtheorem{definition}{Definition}
\title{Flow Proof Artifact: Ring-mul-distrib-right}
\author{Generated by \texttt{flow doc proof}}
\date{Flow Proof Book}
\begin{document}
\maketitle
Ring multiplication distributes over addition on the right.\\[0.5em]
\textbf{Source.} dummit-foote — *Abstract Algebra*, §7.1\\[0.5em]
\bigskip
\noindent{\large \textbf{Derived fact 1.} \textit{(b + c) * a = b * a + c * a} \hfill Ring \cdot multiplication \cdot right distribution over addition holds}\par
\smallskip\noindent\textit{Coordinate: Ring \cdot multiplication \cdot right distribution over addition holds}\par
\smallskip\noindent\textit{Source: dummit-foote}\par
\smallskip\noindent\textit{Built on: left distribution over addition holds, for multiplication on Ring, order does not matter, for addition on Ring}\par
\medskip
\noindent\textbf{Goal.} (b + c) * a = b * a + c * a\par
\noindent$\displaystyle \forall a \in Ring \forall b \in Ring \forall c \in Ring\quad (b + c) * a = b \cdot a + c \cdot a$\par
\medskip
\noindent
\begin{tabular}{@{} >{\raggedright\arraybackslash}p{0.06\textwidth} >{\raggedright\arraybackslash}p{0.47\textwidth} >{\raggedleft\arraybackslash}p{0.41\textwidth} @{}}
\textbf{\#} & \textbf{Proof} & \textbf{Mathematics} \\
\midrule
\textbf{1.} & We prove that right distribution over addition holds for multiplication on Ring. & \\[0.45em]
\textbf{2.} & We invoke the definitional clause governing multiplication on Ring: left distribution over addition holds, for multiplication on Ring (instantiated for a, b, c). & \\[0.45em]
\textbf{3.} & We invoke the definitional clause governing addition on Ring: order does not matter, for addition on Ring (instantiated for b, c). & \\[0.45em]
\textbf{4.} & From step 2 and step 3, this implies (b plus c) times a equals b times a plus c times a. Hence proven. & $\displaystyle (b + c) * a = b \cdot a + c \cdot a$ \\[0.45em]
\bottomrule
\end{tabular}
\label{thm:Ring-multiplication-right_distribution_over_addition_holds}
\par\medskip\hrule
\end{document}
